% BHU PYQ -- Transcribed & Corrected
% Paper: MATMJ-31 : Linear Algebra (Mid-Term)
% Exam: B.Sc. (Hons.) Semester III Mid-Term Examination, 2025-26 (NEP)
% Department: Department of Mathematics

\documentclass[11pt,a4paper]{article}
\usepackage[a4paper,top=2.5cm,bottom=2.5cm,left=2.8cm,right=2.8cm,headheight=14pt]{geometry}
\usepackage{amsmath,amssymb}
\usepackage[T1]{fontenc}
\usepackage[utf8]{inputenc}
\usepackage{lmodern,microtype,fancyhdr,enumitem,hyperref,lastpage,parskip}

\hypersetup{
    colorlinks=true,
    urlcolor=black,
    linkcolor=black,
    pdftitle={MATMJ31: Linear Algebra Mid-Term -- BHU B.Sc. Sem III 2025-26 (NEP)}
}

\pagestyle{fancy}
\fancyhf{}
\renewcommand{\headrulewidth}{0.4pt}
\renewcommand{\footrulewidth}{0.4pt}
\fancyhead[L]{\small   \textit{Banaras Hindu University}}
\fancyhead[R]{\small   \textit{MATMJ-31: Linear Algebra (Mid Term)}}
\fancyfoot[C]{\small Page  \thepage\ of \pageref{LastPage}}


\newlist{parts}{enumerate}{2}
\setlist[parts,1]{label=(\roman*),leftmargin=*,itemsep=4pt,topsep=2pt}

\newcommand{\pts}[1]{\hfill   \textit{[#1]}}


\begin{document}


\begin{center}
  {\large\bfseries BANARAS HINDU UNIVERSITY}\\[2pt]
  {
\normalsize B.Sc. (Hons.) Semester III Mid-Term Examination, 2025-26 (NEP)}\\[6pt]
  \rule{0.8\linewidth}{0.5pt}\\[6pt]
  {\large\bfseries DEPARTMENT OF MATHEMATICS}\\[4pt]
  {\large\bfseries Paper Code: MATMJ-31 : Linear Algebra (Mid-Term)}\\[4pt]
  {
\normalsize Time: 1 Hour\hfill Maximum Marks: 20}
\end{center}

\rule{\linewidth}{0.4pt}

\smallskip



\noindent   \textit{Instructions: Attempt any FIVE questions. All questions carry equal marks (4 marks each).}

\bigskip




\noindent   \textbf{Question 1.} Define a Vector Space over a field $F$. Prove that a non-empty subset $W \subseteq V$ is a subspace of $V$ if and only if $\alpha u + \beta v \in W$ for all $\alpha, \beta \in F$ and $u, v \in W$.\pts{4}

\medskip\rule{\linewidth}{0.2pt}\medskip




\noindent   \textbf{Question 2.} Test the linear independence of vectors $v_1 = (1, 2, 1)$, $v_2 = (2, 1, 0)$, and $v_3 = (1, -1, 2)$ in $\mathbb{R}^3$.\pts{4}

\medskip\rule{\linewidth}{0.2pt}\medskip




\noindent   \textbf{Question 3.} State and prove the Rank-Nullity Theorem for a linear transformation $T \colon V  o W$.\pts{4}

\medskip\rule{\linewidth}{0.2pt}\medskip




\noindent   \textbf{Question 4.} Find the basis and dimension of the subspace $W = \{(x, y, z) \in \mathbb{R}^3 \mid x + 2y - z = 0\}$.\pts{4}

\medskip\rule{\linewidth}{0.2pt}\medskip




\noindent   \textbf{Question 5.} Let $T \colon \mathbb{R}^2  o \mathbb{R}^2$ be defined by $T(x, y) = (2x + y, x - 3y)$. Find the matrix of $T$ relative to the standard basis of $\mathbb{R}^2$.\pts{4}

\vfill
\rule{\linewidth}{0.4pt}

\begin{center}\small
\href{https://bhu-exam.vercel.app/}{Click here to visit our website}\\[4pt]
\href{https://me-aryan.vercel.app/}{Click here to visit our contributor}
\end{center}

\end{document}
